\section{Sistemas legados e modernização}

Um sistema legado é um sistema antigo que ainda está em uso. Ele pode estar escrito em uma linguagem pouco usada hoje, ter uma estrutura difícil de mudar ou depender de bibliotecas sem suporte. Mesmo assim, ele costuma ser importante, porque guarda regras de negócio que a organização usa todos os dias.

Um exemplo é o sistema estudado por \citeonline{alves2026legacy}. Era um sistema de gestão em Visual Basic 6, em produção havia mais de vinte anos. Outro exemplo são os programas em COBOL, que a IBM migra para Java \cite{chakaravarthy2026cobol}.

A modernização é o trabalho de levar esse sistema para uma tecnologia mais nova, sem perder o que ele já faz. Ela pode ser de dois tipos:

\begin{itemize}
    \item \textbf{conversão de linguagem}, por exemplo de Visual Basic 6 para C\#;
    \item \textbf{atualização de versão}, por exemplo de Java 8 para Java 17 \cite{liu2026migrationbench, may2025freshbrew}.
\end{itemize}

Uma dificuldade comum é a falta de documentação. Muitas vezes, o código é a única fonte do que o sistema faz. \citeonline{diggs2024leveraging} testaram modelos de linguagem para documentar código legado em MUMPS e em assembly. Entender o código antigo é o primeiro passo da modernização.
